\chapter*{Chapter 8 추가 기호와 용어}
\addcontentsline{toc}{chapter}{Chapter 8 추가 기호와 용어}
\markboth{Chapter 8 추가 기호와 용어}{Chapter 8 추가 기호와 용어}

\renewcommand{\arraystretch}{1.2}
\begin{longtable}{>{\raggedright\arraybackslash}p{0.27\textwidth}>{\raggedright\arraybackslash}p{0.63\textwidth}}
\toprule
\textbf{기호·용어} & \textbf{뜻} \\
\midrule
\endfirsthead
\toprule
\textbf{기호·용어} & \textbf{뜻} \\
\midrule
\endhead
\bottomrule
\endfoot

기본 작도 단계 & 두 직선, 직선과 원, 또는 두 원의 교점을 새 점으로 추가하는 조작. \\
$\mathcal C(M)$ & 시작점 집합 $M\subseteq\C$에서 자와 컴퍼스로 유한번의 작도를 거쳐 얻는 모든 점의 집합. \\
$\mathcal C_0$ & $\mathcal C(\{0,1\})$. 표준 작도가능체. \\
$\mathcal C_{\R}$ & $\mathcal C_0\cap\R$. 실수 작도가능체. \\
작도가능점·작도가능수 & $\mathcal C(M)$ 또는 $\mathcal C_0$에 속하는 점·수. \\
이차확장 탑 & 각 단계의 차수가 $2$ 이하인 체확장 사슬 $K=L_0\subseteq\cdots\subseteq L_r$. \\
갈루아 $2$-확장 & 갈루아군이 유한 $2$-군인 유한 갈루아 확장. 확장차수가 $2$의 거듭제곱이다. \\
정상폐포 판정 & 대수적 수가 작도가능할 필요충분조건이 그 최소다항식의 분해체 갈루아군이 $2$-군인 것이라는 판정. \\
정육면체 배적 & 주어진 정육면체 부피의 두 배인 정육면체를 작도하는 고전 문제. 단위 경우 목표 길이는 $\sqrt[3]2$. \\
각의 삼등분 & 임의의 주어진 각을 자와 컴퍼스로 세 등분하는 고전 문제. $60^\circ$가 반례이다. \\
원적문제 & 주어진 원과 같은 넓이의 정사각형을 작도하는 문제. 단위원에서는 목표 변 길이가 $\sqrt\pi$이다. \\
정$n$각형 작도 판정 & 정$n$각형이 작도가능할 필요충분조건은 $\varphi(n)$이 $2$의 거듭제곱인 것이다. \\
페르마 수 & $F_r=2^{2^r}+1$. \\
페르마 소수 & 소수인 페르마 수. \\
가우스--방첼 정리 & 정$n$각형이 작도가능할 필요충분조건은 $n=2^m p_1\cdots p_r$이고 $p_i$들이 서로 다른 페르마 소수인 것이다. \\
황금비 & $\phi=(1+\sqrt5)/2$. 정오각형의 대각선과 변의 비. \\

\end{longtable}
